\usepackage[scr=rsfs,cal=euler]{mathalfa}
\let\mathds\mathbb
\let\epsilon\varepsilon
\newcommand{\Changel}{\mathcode`l="7160}
\newcommand{\Changelback}{\mathcode`l="716C}
\newcommand{\NoChangel}[1]{%
\expandafter\let\csname old\string#1\endcsname=#1
\let#1=\relax
\newcommand{#1}{\mathcode`l="716C\csname old\string#1\endcsname\mathcode`l="7160 }%
}

\NoChangel{\log}
\NoChangel{\ln}
\NoChangel{\lim}
\NoChangel{\limsup}
\NoChangel{\liminf}
\NoChangel{\varlimsup}
\NoChangel{\varliminf}
\NoChangel{\varinjlim}
\NoChangel{\varprojlim}

\newenvironment{enumeratei}
{\bgroup\def\theenumi{\roman{enumi}}\def\theenumii{\arabic{enumii}}\begin{enumerate}}
{\end{enumerate}\egroup}
\newcommand{\Psfrac}[2]{(\sfrac{#1}{#2})}
\newcommand{\psfrac}[2]{\sfrac{(#1)}{#2}}
\newcommand{\spfrac}[2]{\sfrac{#1}{(#2)}}
\newcommand\mto{\mathchoice{\longmapsto}{\mapsto}{\mapsto}{\mapsto}}
\newcommand\hto{\mathrel{\lhook\joinrel\to}}

\usepackage{graphicx}

\newtheorem{thm}{Theorem}[section]
\newtheorem{cor}[thm]{Corollary}
\newtheorem{lem}[thm]{Lemma}
\newtheorem{prop}[thm]{Proposition}
\newtheorem{quest}{Question}[section]
\newtheorem{conj}{Conjecture}[section]
\newtheorem{maintheorem}{Theorem}
\theoremstyle{definition}
\newtheorem{rem}[thm]{Remark}
\newtheorem{exm}[thm]{Example}
\newtheorem{ass}[thm]{Assumption}
\renewcommand*{\themaintheorem}{\Alph{maintheorem}}

\DeclareMathOperator{\area}{area}
\DeclareMathOperator{\lcm}{lcm}
\DeclareMathOperator{\const}{const}
\DeclareMathOperator{\Cal}{Cal}
\DeclareMathOperator{\Flux}{\mathrm{Flux}}
\newcommand{\id}{\mathrm{id}}
\newcommand{\Ham}{\mathrm{Ham}}
\newcommand{\std}{\mathrm{std}}
\newcommand{\vol}{\mathrm{vol}}
\newcommand{\fix}{\mathrm{fix}}
\newcommand{\pr}{\mathrm{pr}}
\newcommand{\Tan}{\mathrm{T}}
\newcommand{\N}{\mathds{N}}
\newcommand{\Z}{\mathds{Z}}
\newcommand{\R}{\mathds{R}}
\newcommand{\CC}{\mathds{C}}\let\C\CC
\newcommand{\UU}{\mathcal{U}}
\newcommand{\VV}{\mathcal{V}}

\let\eul e
\let\diff d
